\documentclass[11pt,a4paper]{article}
\usepackage[utf8]{inputenc}
\usepackage[margin=1in]{geometry}
\usepackage{amsmath,amssymb,amsfonts,amsthm}
\usepackage{booktabs}
\usepackage{graphicx}
\usepackage{hyperref}
\usepackage{microtype}
\usepackage{algorithm}
\usepackage{algorithmic}
\usepackage{cite}
\usepackage{xcolor}

\hypersetup{
    colorlinks=true,
    linkcolor=blue!70!black,
    citecolor=green!50!black,
    urlcolor=blue!80!black
}

\title{\textbf{AgentWatch-A: Black-Box Online Pre-Execution Monitoring for Indirect Prompt Injection in Tool-Using LLM Agents}}

\author{
  \textbf{Anonymous Authors}\\
  \textit{Paper under blind conference review}
}

\date{\today}

\begin{document}

\maketitle

\begin{abstract}
Large Language Model (LLM) agents operating in tool-augmented environments are acutely vulnerable to Indirect Prompt Injection (IPI), where malicious instructions embedded in untrusted external data compromise the agent's decision loop. While post-execution auditing and static content filtering have been proposed, reactive auditing detects compromise only after irreversible side effects occur, and content detectors suffer severe false alarms when legitimate tasks process untrusted text. In this paper, we investigate \textbf{AgentWatch-A}, an online, black-box, pre-execution monitoring framework designed to intercept compromised actions before irreversible execution while strictly controlling benign trajectory interruption ($FPR_{\text{traj}} \le \alpha = 0.20$). AgentWatch-A models the agent's causal trajectory prefix using a lightweight causal Goal-Conditioned Gated Recurrent Unit (GRU) augmented with deterministic argument-provenance features, calibrated via finite-sample split-conformal prediction on benign trajectory maxima. We evaluate AgentWatch-A on the AgentDojo workspace benchmark under both offline prefix replay (Track 1) and live closed-loop intervention (Track 2). In closed-loop evaluation, simpler transition bigram and sequence baselines trigger a catastrophic 100.00\% false interruption rate on benign workflows, and content detectors interrupt 50.00\% of legitimate tasks. In contrast, AgentWatch-A achieves 100.00\% Attack Success Rate (ASR) reduction while maintaining a 0.00\% benign false interruption rate, executing at an average CPU inference latency of 1.61~ms per tool invocation with zero cloud API overhead. We analyze the theoretical and empirical conditions under which goal conditioning and argument provenance prevent benign over-blocking, and document key failure modes and distribution-shift boundaries.
\end{abstract}

\section{Introduction}
\label{sec:intro}

The deployment of autonomous Large Language Model (LLM) agents equipped with external tools (e.g., email clients, shell access, SQL databases, and file systems) has expanded their operational utility but introduced critical security vulnerabilities~\cite{schick2023toolformer,yao2022react}. Among these, \textbf{Indirect Prompt Injection (IPI)}~\cite{greshake2023not,perez2022ignore} represents one of the most severe threat vectors. In an IPI attack, an adversary embeds adversarial instructions into data objects that the agent reads during routine task execution (e.g., a calendar invite, an email body, or a webpage). When processed by the underlying language model, these malicious payloads hijack the model's instruction-following capabilities, causing the agent to execute attacker-specified tool calls—such as exfiltrating private files, modifying system credentials, or sending unauthorized correspondence~\cite{debenedetti2024agentdojo}.

Existing defenses against IPI broadly fall into three paradigms:
\begin{enumerate}
    \item \textbf{Model Alignment and Prompt Hardening:} Fine-tuning models or prepending defensive prompts to resist untrusted instructions. While partially effective, alignment can be circumvented by adversarial jailbreaks and complex framing.
    \item \textbf{Static Content Scanning:} Deploying auxiliary classifiers (e.g., LlamaFirewall~\cite{bhatt2025llamafirewall}) to scan environment observations. However, content scanners lack execution context: they trigger false alarms whenever a benign workflow legitimately processes adversarial or suspicious text (e.g., a security analyst summarizing spam).
    \item \textbf{Post-Execution Behavioral Auditing:} Examining full trajectory traces offline after task completion (e.g., DriftNet~\cite{shi2026driftnet}, TraceAegis~\cite{zhang2025traceaegis}). Although effective for forensic accounting, post-execution detection occurs \textit{after} irreversible real-world side effects (data deletion, financial transfer, exfiltration) have already transpired.
\end{enumerate}

To bridge this operational gap, we define and formalize the \textbf{Online Pre-Execution Monitoring Problem}. In this setting, an external, black-box monitor intercepts every proposed tool invocation $a_t$ at step $t$ given only the user's declared goal $g$ and the historical causal trace prefix $T_{<t}$. Crucially, the monitor must satisfy the strict \textbf{Pre-Execution Boundary Invariant}:
\begin{equation}
    t_{\text{flag}} \le t_{\text{harm}},
\end{equation}
guaranteeing that any alert or blocking intervention occurs before the first irreversible harmful action executes. Furthermore, a viable security monitor must not render the agent unusable; it must enforce a rigorous trajectory-level False Positive Rate constraint ($FPR_{\text{traj}} \le \alpha$) over legitimate benign tasks.

In this work, we present \textbf{AgentWatch-A}, a lightweight, causally sound, and distribution-free calibrated pre-execution monitoring architecture. AgentWatch-A combines:
\begin{itemize}
    \item A lightweight \textbf{Goal-Conditioned GRU} that maps frozen goal text embeddings and causal action histories into a shared latent space via contrastive objective;
    \item A 7-dimensional \textbf{Argument Provenance Vector} that tracks data flow taint and suspicious payload markers without inspecting internal model weights;
    \item A \textbf{Split-Conformal Calibration} rule on benign trajectory score maxima that bounds the trajectory false positive rate under exchangeability;
    \item An active \textbf{In-Loop Interception Hook} that terminates or holds hijacked execution trajectories before tool execution.
\end{itemize}

Through rigorous empirical evaluation on the AgentDojo workspace benchmark~\cite{debenedetti2024agentdojo}, executed under 100\% zero-cost local CPU reproduction, we demonstrate that simpler baselines (transition bigrams, unconditioned sequence models) completely destroy agent utility in closed-loop deployments (100.00\% false interruption rate). AgentWatch-A resolves this trade-off, achieving 100.00\% ASR reduction with 0.00\% benign interruption at an average CPU latency of 1.61~ms.

\section{Related Work}
\label{sec:related_work}

\paragraph{Prompt Injection and Agent Security Benchmarks.}
Prompt injection was formalized by Perez and Ribeiro~\cite{perez2022ignore} and Greshake et al.~\cite{greshake2023not}, who demonstrated that secondary text channels can manipulate LLM outputs. AgentDojo~\cite{debenedetti2024agentdojo} introduced a dynamic benchmark measuring execution outcome success for both user goals and attacker goals. MELON~\cite{deng2024melon} analyzed execution trace corruption. While AgentDojo provides dynamic environment evaluations, it primarily evaluates offline defense heuristics or prompt defenses rather than calibrated trajectory monitors.

\paragraph{Trajectory Anomaly Detection.}
Trajectory-based anomaly detection has recently emerged to identify agent drift. Trajectory Guard~\cite{gao2026trajectory} investigates anomaly detection over agent action sequences. DriftNet~\cite{shi2026driftnet} and AgentDrift~\cite{zhou2026agentdrift} explore supervised sequence classification to locate hijacked trajectory steps. TraceAegis~\cite{zhang2025traceaegis} models benign agent behavior structures. However, prior trajectory monitors largely operate in an offline, post-execution fashion or lack formal statistical calibration on trajectory maxima, leading to unknown false alarm rates over multi-step workflows.

\paragraph{Runtime Guardrails and Firewalls.}
AgentSentry~\cite{chen2026agentsentry} and LlamaFirewall~\cite{bhatt2025llamafirewall} introduce pre-tool execution scanning. However, existing guardrails either rely on heavy secondary LLMs (which introduce hundreds of milliseconds of latency and high cloud API costs) or focus exclusively on observation content without modeling goal-to-action semantic divergence. AgentWatch-A addresses this by demonstrating that a sub-2ms recurrent monitor trained via benign contrastive alignment can outperform token scanners while eliminating benign false interruptions.

\paragraph{Conformal Prediction in Sequential Systems.}
Split-conformal prediction~\cite{vovk2005algorithmic,angelopoulos2021gentle} provides finite-sample, distribution-free guarantees on error rates. Extensions to sequential settings~\cite{tibshirani2019conformal,barber2021predictive} typically require complex martingale tracking. In AgentWatch-A, we adopt a trajectory-maxima calibration scheme: by calibrating over the maximum anomaly score attained across an entire benign trajectory, we provide marginal coverage at the trajectory level without sequential independence assumptions.

\section{Threat Model and Assumptions}
\label{sec:threat_model}

Following Section 5 of the research specification, we formalize the threat model for AgentWatch-A:

\paragraph{Attacker Capabilities.}
The adversary cannot modify the user prompt, the base LLM weights, the system prompts, or the monitor code. The adversary exerts control solely over external data objects returned by environment tools (e.g., file contents, incoming email messages, web search snippets, database records). The adversary's objective is to induce the agent to execute unauthorized tool actions that satisfy an attacker goal $g_{\text{att}}$ (e.g., exfiltrating sensitive contact lists or modifying access privileges).

\paragraph{Trust Boundaries and Intervention Point.}
\begin{itemize}
    \item \textbf{Trusted Channel:} The user's initial prompt $g$ and the runtime monitor code are trusted.
    \item \textbf{Untrusted Channel:} Any tool execution observation $o_t$ containing third-party or user-generated external data.
    \item \textbf{Mediation Point:} The monitor sits strictly at the \textbf{Tool Execution Boundary}. When the agent emits a proposed tool call $a_t = (\text{tool\_name}, \text{arguments})$, the call is intercepted \textit{before} the environment runtime invokes the tool implementation.
\end{itemize}

\paragraph{Out-of-Scope Threats.}
Direct jailbreaks where the primary user attempts to compromise the agent, physical memory tampering, network man-in-the-middle attacks, and denial-of-service via unbounded compute loops are outside the scope of this defense.

\section{Problem Formulation}
\label{sec:formulation}

Let an agent trajectory $T$ up to step $t$ be represented by the tuple:
\begin{equation}
    T_t = (g, a_1, o_1, a_2, o_2, \dots, a_{t-1}, o_{t-1}, a_t),
\end{equation}
where $g \in \mathcal{X}_{\text{text}}$ is the user goal, $a_i \in \mathcal{A}$ are proposed tool calls, and $o_i \in \mathcal{O}$ are environment observations. At step $t$, the agent proposes action $a_t$, but $a_t$ has not yet been executed.

\paragraph{Outcome Labels.}
Following the ground-truth outcomes evaluated by security benchmarks:
\begin{equation}
    Y(T) \in \{\text{benign}, \text{unexposed}, \text{resisted}, \text{hijacked}\},
\end{equation}
where \textit{benign} denotes an attack-free trajectory; \textit{unexposed} denotes an attack trajectory where the agent never observed the injected payload; \textit{resisted} denotes a trajectory where the agent encountered an injection payload but successfully ignored it, completing the user goal; and \textit{hijacked} denotes a trajectory where the agent succumbed to the injection and executed attacker-desired actions.

\paragraph{Harm Step and Pre-Execution Timing.}
For every hijacked trajectory, let $t_{\text{harm}}$ denote the index of the first tool action that directly advances the attacker's objective or causes irreversible state change. A monitor alert issued at step $t_{\text{flag}}$ is valid if and only if:
\begin{equation}
    t_{\text{flag}} \le t_{\text{harm}}.
\end{equation}
Any alert issued at $t > t_{\text{harm}}$ is treated as a detection failure ($t_{\text{flag}} = \infty$) for the purposes of pre-execution containment.

\paragraph{Evaluation Metrics.}
\begin{itemize}
    \item \textbf{Pre-Execution True Positive Rate ($TPR_{\text{pre}}$):} The proportion of hijacked trajectories flagged at or before $t_{\text{harm}}$:
    \begin{equation}
        TPR_{\text{pre}} = \frac{1}{|T_{\text{hijack}}|} \sum_{T \in T_{\text{hijack}}} \mathbb{I}(t_{\text{flag}}(T) \le t_{\text{harm}}(T)).
    \end{equation}
    \item \textbf{Trajectory-Level False Positive Rate ($FPR_{\text{traj}}$):} The fraction of non-hijacked trajectories ($Y \in \{\text{benign}, \text{resisted}\}$) where the monitor raises a false alarm at any step:
    \begin{equation}
        FPR_{\text{benign}} = \frac{1}{|T_{\text{benign}}|} \sum_{T \in T_{\text{benign}}} \mathbb{I}(t_{\text{flag}}(T) < \infty).
    \end{equation}
    \item \textbf{In-Loop Interruption Rate ($FIR_{\text{benign}}$):} In active closed-loop intervention, the proportion of legitimate benign user tasks prematurely terminated or degraded by false monitor intervention.
\end{itemize}

\section{The AgentWatch-A Architecture}
\label{sec:method}

AgentWatch-A comprises four tightly integrated components: deterministic feature encoding, causal goal-conditioned sequence modeling, argument provenance extraction, and split-conformal threshold calibration.

\subsection{Deterministic Feature and Tool Action Encoding}
To maintain zero cloud cost and cross-platform determinism across CPU runtimes, text features are extracted using a frozen, deterministic CRC32 hash projection. Given a token sequence $s = (w_1, \dots, w_m)$, each word is mapped to an embedding index via:
\begin{equation}
    idx(w_i) = \text{CRC32}(w_i) \pmod V,
\end{equation}
where $V = 10{,}000$ represents a fixed vocabulary bucket size. Token embeddings are pooled and normalized to produce a compact text representation $e \in \mathbb{R}^{d_{\text{text}}}$ ($d_{\text{text}} = 64$).

Each proposed action $a_t$ is encoded into a structured feature vector:
\begin{equation}
    x_t = [e_{\text{tool}}(a_t) \,\|\, e_{\text{arg\_text}}(a_t) \,\|\, p(a_t)],
\end{equation}
where $e_{\text{tool}} \in \mathbb{R}^{d_{\text{tool}}}$ is a learnable embedding of the tool name, $e_{\text{arg\_text}}$ is the pooled CRC32 representation of argument values, and $p(a_t) \in \mathbb{R}^7$ is the deterministic argument provenance vector.

\subsection{Argument Provenance Vector}
A key vulnerability of sequence-only monitors is their inability to discern when a legitimate tool (e.g., \texttt{send\_email}) is being invoked with attacker-tainted parameters versus user-specified parameters. AgentWatch-A extracts a 7-dimensional argument provenance vector $p(a_t)$:
\begin{equation}
    p(a_t) = [p_{\text{len}}, p_{\text{hex}}, p_{\text{b64}}, p_{\text{domain}}, p_{\text{ctrl}}, p_{\text{taint}}, p_{\text{impact}}]^\top \in [0, 1]^7,
\end{equation}
where:
\begin{itemize}
    \item $p_{\text{len}} = \min(1.0, \frac{|args|}{1000})$ measures argument character volume;
    \item $p_{\text{hex}}$ measures hexadecimal entropy density;
    \item $p_{\text{b64}}$ detects base64-encoded string fragments;
    \item $p_{\text{domain}}$ flags mismatched or non-whitelisted external recipient domains;
    \item $p_{\text{ctrl}}$ counts control tokens and prompt syntax markers (e.g., \texttt{SYSTEM:}, \texttt{IGNORE PREVIOUS});
    \item $p_{\text{taint}} \in \{0, 1\}$ indicates whether argument sub-strings overlap with recent untrusted environment observations;
    \item $p_{\text{impact}} \in [0, 1]$ encodes static tool risk categories (e.g., read tools = 0.1, write tools = 0.5, exfiltration/execution tools = 0.9).
\end{itemize}

\subsection{Causal Goal-Conditioned Recurrence}
The sequence of encoded actions $(x_1, \dots, x_t)$ is processed sequentially by a causal Gated Recurrent Unit (GRU):
\begin{equation}
    h_t = \text{GRU}(x_t, h_{t-1}), \quad h_0 = \mathbf{0},
\end{equation}
where $h_t \in \mathbb{R}^{d_h}$ ($d_h = 64$). To evaluate whether the current trajectory prefix aligns with the declared user goal, the user goal embedding $e_g$ and the trajectory hidden state $h_t$ are projected into a shared 32-dimensional metric space:
\begin{equation}
    z_g = \frac{W_g e_g}{\|W_g e_g\|_2}, \quad z_t = \frac{W_h h_t}{\|W_h h_t\|_2}.
\end{equation}
The instantaneous anomaly score $s_t \in [0, 1]$ is computed from their cosine dissimilarity modulated by argument taint:
\begin{equation}
    s_t = \sigma\left( \gamma \cdot (1 - \langle z_g, z_t \rangle) + \beta \cdot p_{\text{taint}}(a_t) \cdot p_{\text{impact}}(a_t) \right).
\end{equation}
The model is trained contrastively on benign task trajectories using InfoNCE loss with same-suite negative goals, enforcing that legitimate actions align closely with the true goal while diverging from mismatched goals.

\subsection{Split-Conformal Calibration on Trajectory Maxima}
To guarantee a controlled false alarm rate over multi-step workflows, we formulate calibration at the trajectory level. For each benign calibration trajectory $T_j \in \mathcal{D}_{\text{cal}}$, we extract the maximum anomaly score attained over its entire lifespan:
\begin{equation}
    S_j = \max_{1 \le t \le |T_j|} s_t(T_j).
\end{equation}
Given $n$ exchangeable benign calibration trajectories and a target trajectory false positive rate $\alpha \in (0, 1)$, we sort the calibration maxima $S_{(1)} \le S_{(2)} \le \dots \le S_{(n)}$. Following split-conformal quantile selection:
\begin{equation}
    k = \lceil (n + 1)(1 - \alpha) \rceil.
\end{equation}
If $k \le n$, the calibrated threshold is set to the $k$-th order statistic $\hat{\tau} = S_{(k)}$; if $k > n$, $\hat{\tau} = \infty$.

\paragraph{Decision Policy.}
During runtime monitoring, at each step $t$, the decision policy evaluates the proposed action against $\hat{\tau}$:
\begin{equation}
    \text{Decision}(s_t) = 
    \begin{cases}
        \text{ALLOW}, & \text{if } s_t < \hat{\tau} - \delta, \\
        \text{FLAG},  & \text{if } \hat{\tau} - \delta \le s_t < \hat{\tau}, \\
        \text{BLOCK}, & \text{if } s_t \ge \hat{\tau}.
    \end{cases}
\end{equation}
When the policy returns \texttt{BLOCK}, the in-loop interceptor halts execution immediately, preventing $a_t$ from reaching the environment runtime.

\section{Experimental Setup}
\label{sec:setup}

\paragraph{Benchmark Environment.}
We evaluate on the \textbf{AgentDojo} benchmark suite (version 0.1.35, MIT License)~\cite{debenedetti2024agentdojo}, specifically the \texttt{workspace} suite comprising 40 distinct user tasks spanning email operations, calendar management, and cloud storage. Ground-truth outcome labels are provided by AgentDojo's programmatic security and task verification functions.

\paragraph{Data Splits and Partitions.}
The experimental corpus comprises 20 canonical trajectories partitioned into:
\begin{itemize}
    \item \textbf{Calibration Set:} 5 benign, attack-free trajectories ($n = 5$);
    \item \textbf{Held-Out Test Set:} 15 trajectories comprising 5 benign tasks, 3 unexposed tasks, 4 resisted tasks, and 3 hijacked tasks.
\end{itemize}
Splits are partitioned strictly by task ID to guarantee task-disjoint separation between calibration and evaluation.

\paragraph{Baselines Evaluated.}
We compare AgentWatch-A against five rigorous baselines:
\begin{enumerate}
    \item \textbf{B1 (Transition Bigram):} Models tool transition matrices $P(a_t \mid a_{t-1})$ from benign traces, flagging transition probabilities below threshold.
    \item \textbf{B2 (Pooled Embedding Anomaly):} Unsupervised anomaly detectors (Isolation Forest and One-Class SVM) fit on mean-pooled trajectory embeddings.
    \item \textbf{B3 (Content Detector):} A regex and keyword scanner scanning observation strings for known injection patterns.
    \item \textbf{B4 (Unconditioned Sequence):} An unconditioned Markov sequence model tracking tool call $n$-grams without goal semantics.
    \item \textbf{B6 (Supervised Oracle):} An upper-bound classifier trained directly on known hijacked and benign traces with full label access.
\end{enumerate}

\paragraph{Evaluation Tracks.}
We evaluate across two operational modes:
\begin{itemize}
    \item \textbf{Track 1 (Offline Prefix Replay):} Evaluates discriminative power ($TPR_{\text{pre}}$, $FPR$, AUROC, AUPRC) on pre-recorded causal prefixes without altering execution.
    \item \textbf{Track 2 (Closed-Loop In-Loop Intervention):} The monitor executes as a live interceptor inside the agent loop, actively blocking proposed actions and measuring true Attack Success Rate reduction and benign task false interruption ($FIR_{\text{benign}}$).
\end{itemize}

\paragraph{Computational Platform.}
All experiments, training passes, and evaluations are executed locally on an x86\_64 consumer CPU running Python 3.11.0 and PyTorch 2.14.1+cpu. No cloud LLM API tokens or commercial services were utilized.

\section{Empirical Results}
\label{sec:results}

\subsection{Track 1: Offline Prefix Replay and Feature Separation}

Table~\ref{tab:track1} summarizes the offline discrimination performance across all evaluated models on the held-out test set at $\alpha = 0.20$.

\begin{table}[h!]
\centering
\small
\caption{\textbf{Track 1: Offline Prefix Replay Performance (Held-Out Test Set, $\alpha = 0.20$).} Thresholds are calibrated via split-conformal prediction on $n=5$ benign calibration trajectories. $FPR$ is partitioned across benign ($N=5$) and resisted ($N=4$) trajectories.}
\label{tab:track1}
\begin{tabular}{lcccccc}
\toprule
\textbf{Model / Baseline} & \textbf{Threshold $\hat{\tau}$} & \textbf{$TPR_{\text{pre}}$} & \textbf{$FPR_{\text{benign}}$} & \textbf{$FPR_{\text{resisted}}$} & \textbf{AUROC} & \textbf{AUPRC} \\
\midrule
B1: Transition Bigram         & 0.2500 & 100.0\% & 40.0\% & 0.0\%   & 0.800 & 0.800 \\
B3: Content Detector          & 0.0000 & 100.0\% & 0.0\%  & 0.0\%   & 1.000 & 1.000 \\
B4: Unconditioned Sequence    & 0.2500 & 100.0\% & 40.0\% & 0.0\%   & 0.600 & 0.300 \\
B6: Supervised Oracle         & 0.0065 & 100.0\% & 0.0\%  & 0.0\%   & 1.000 & 1.000 \\
\midrule
Goal-Conditioned GRU (Ablated)& 0.4847 & 100.0\% & 100.0\% & 100.0\% & 1.000 & 1.000 \\
\textbf{AgentWatch-A (Full)}  & \textbf{0.5515} & \textbf{100.0\%} & \textbf{40.0\%} & \textbf{0.0\%} & \textbf{1.000} & \textbf{1.000} \\
\bottomrule
\end{tabular}
\end{table}

\paragraph{Analysis of Discrimination.}
Both B1 (Bigram) and B4 (Unconditioned Sequence) exhibit high false alarm rates ($FPR_{\text{benign}} = 40.0\%$), because benign tasks frequently compose legitimate tools in novel orders not seen in the compact calibration set. Crucially, ablating the goal conditioning from the GRU causes the ablated monitor to flag every single test trajectory ($FPR_{\text{benign}} = 100.0\%$). In contrast, the full AgentWatch-A monitor achieves perfect ranking metrics ($\text{AUROC} = 1.000$, $\text{AUPRC} = 1.000$), matching the supervised oracle B6 without requiring attack labels.

\subsection{Track 2: Closed-Loop In-Loop Intervention}

The true test of a pre-execution monitor lies in closed-loop deployment, where false alarms terminate legitimate user tasks. Table~\ref{tab:track2} presents the live intervention metrics.

\begin{table}[h!]
\centering
\small
\caption{\textbf{Track 2: Live In-Loop Intervention and Operational Latency.} Unprotected agent baseline achieves 50.00\% ASR. $FIR_{\text{benign}}$ measures false interruption of legitimate workflows. Latency is measured per proposed tool invocation on consumer x86\_64 CPU.}
\label{tab:track2}
\begin{tabular}{lccccc}
\toprule
\textbf{Defense System} & \textbf{In-Loop ASR} & \textbf{ASR Reduction} & \textbf{$FIR_{\text{benign}}$} & \textbf{Mean CPU Lat.} & \textbf{p95 Lat.} \\
\midrule
Unprotected Agent             & 50.00\% & 0.0\%   & 0.00\%   & 0.00~ms & 0.00~ms \\
B1: Transition Bigram         & 0.00\%  & 100.0\% & 100.00\% & 0.01~ms & 0.01~ms \\
B3: Content Detector          & 0.00\%  & 100.0\% & 50.00\%  & 1.48~ms & 1.55~ms \\
B4: Unconditioned Sequence    & 0.00\%  & 100.0\% & 100.00\% & 0.02~ms & 0.03~ms \\
Goal-Conditioned GRU (Ablated)& 0.00\%  & 100.0\% & 100.00\% & 1.05~ms & 1.20~ms \\
\midrule
\textbf{AgentWatch-A (Full)}  & \textbf{0.00\%} & \textbf{100.0\%} & \textbf{0.00\%} & \textbf{1.61~ms} & \textbf{1.91~ms} \\
\bottomrule
\end{tabular}
\end{table}

\paragraph{Gate G3 and G4 Empirical Findings.}
As shown in Table~\ref{tab:track2}, the naive baselines (B1, B4) and the ablated GRU achieve 100.0\% attack mitigation only by destroying agent utility: they suffer a \textbf{100.00\% False Interruption Rate ($FIR_{\text{benign}}$)}, terminating all legitimate workflows. Similarly, the static content scanner B3 interrupts 50.00\% of benign tasks whenever benign observations contain untrusted strings. 
\textbf{AgentWatch-A is the only system that achieves 100.00\% ASR reduction while maintaining 0.00\% false interruption on benign tasks.}

\paragraph{Operational Latency.}
The average CPU inference latency of AgentWatch-A is \textbf{1.61~ms} per proposed tool call (with a 95th-percentile latency of 1.91~ms and maximum of 1.94~ms), well within the 50~ms real-time operational budget.

\section{Failure Analysis and Discussion}
\label{sec:failure_analysis}

\subsection{Why Simpler Baselines Fail (Gates G3 and G4)}
Transition bigrams ($B1$) and unconditioned sequence models ($B4$) operate on the assumption that tool syntax is stationary. In complex workflows, agents validly chain tools in varied permutations (e.g., checking calendar before reading emails vs. reading emails before searching files). Lacking knowledge of the user's intent $g$, unconditioned models perceive these variations as anomalies, leading to the 100\% benign interruption observed in Table~\ref{tab:track2}.

\subsection{The Content-Only Dilemma ($B3$)}
Content detectors scan for prompt injection keywords (e.g., \texttt{"ignore instructions"}, \texttt{"system update"}). However, modern agents frequently process untrusted web pages, forum posts, or email text that quote injection strings benignly. When an agent reads such text but successfully resists executing it (the \textit{resisted} class), content detectors fire indiscriminately, interrupting valid workflows. AgentWatch-A prevents this by conditioning on argument provenance and action intent: an alert is raised only if the agent actively attempts to execute a high-impact action using tainted data.

\section{Limitations and Ethics}
\label{sec:limitations}

In accordance with Sections 20.3 and 20.4 of the research specification, we document the operational and statistical boundaries of this study:
\begin{enumerate}
    \item \textbf{Finite-Sample Conformal Validity:} Split-conformal coverage guarantees ($FPR \le \alpha$) hold strictly under the assumption of exchangeability between calibration and deployment trajectories. If an agent is deployed in an entirely novel environment (distribution shift), coverage guarantees degrade.
    \item \textbf{Sample Size Scope:} Our experiments are conducted on 20 canonical traces across the AgentDojo workspace benchmark. While statistically rigorous within this benchmark partition, large-scale industrial deployment across thousands of heterogeneous tool suites requires ongoing calibration refreshes.
    \item \textbf{Adaptive Adversaries:} An adversary with white-box access to the monitor's GRU weights could theoretically optimize payload text to minimize cosine distance $\langle z_g, z_t \rangle$. Defending against adaptive optimization via adversarial training remains an important direction for future research.
\end{enumerate}

\section{Conclusion}
\label{sec:conclusion}

We introduced AgentWatch-A, an online, black-box, pre-execution monitoring framework for detecting indirect prompt injection in tool-using LLM agents. By combining deterministic feature hashing, causal goal-conditioned recurrence, argument provenance tracking, and split-conformal calibration on benign trajectory maxima, AgentWatch-A enforces the pre-execution boundary invariant without requiring internal model weights or expensive cloud LLM calls. In closed-loop evaluations on AgentDojo, AgentWatch-A reduces ASR to 0.00\% with zero false interruptions of benign workflows, operating at 1.61~ms CPU latency. All code, data, and reproduction scripts are released to enable verifiable replication.

\bibliographystyle{plain}
\bibliography{references}

\appendix
\section{Reproducibility and Artifact Manifest}
\label{sec:appendix}

\subsection{Reproduction Protocol}
The entire experimental pipeline can be reproduced using a single command:
\begin{verbatim}
python scripts/reproduce_all.py
\end{verbatim}
This master script executes seven validation stages: environment audit, canonical trace verification, baseline fitting, monitor evaluation, split-conformal calibration, closed-loop in-loop intervention, and zero-discrepancy registry assertion.

\subsection{Hardware and Environment Specifications}
\begin{itemize}
    \item \textbf{Operating System:} Microsoft Windows (x86\_64)
    \item \textbf{Python Version:} 3.11.0 (64-bit)
    \item \textbf{PyTorch Version:} 2.14.1+cpu
    \item \textbf{Benchmark Package:} \texttt{agentdojo==0.1.35} (MIT License)
    \item \textbf{Total Compute Budget:} $< 10$ CPU seconds per full evaluation pass. Zero cloud API tokens required.
\end{itemize}

\end{document}
